% !TeX root = ../../report.tex
\subsection{Explainable AI (XAI)}\label{sec:xai}

Explainable AI (XAI) encompasses a suite of machine learning (ML) techniques designed to make the decision-making processes of complex, opaque models understandable to human users \citep{arrieta2020}. As AI systems are increasingly deployed in high-stakes domains such as healthcare and law, understanding their internal mechanisms is imperative for algorithmic accountability and legitimacy \citep{arrieta2020,guidotti2021}.

Although interpretability and explainability are frequently used interchangeably in the literature, they represent distinct conceptual paradigms \citep{arrieta2020}. Arrieta et al.\ \citeyearpar{arrieta2020} define interpretability (ante-hoc) as a passive property of models that are transparent by design, such as linear regression, decision trees, and rule-based systems. Their internal decision logic is naturally understandable to a human observer without requiring external diagnostic tools. Explainability (post-hoc), in contrast, refers to an active procedure applied to trained black-box architectures, such as deep neural networks, to extract surrogate representations of decision-making logic. The effectiveness of a post-hoc explanation depends heavily on the target audience and their cognitive capacity to interpret the output.

Biecek and Samek \citeyearpar{biecek2024} categorise XAI methodologies into two cultures according to the primary intent of explanation \citep{biecek2024}. \textit{Blue XAI} (human-value oriented) focuses on justifying decisions to build trust, satisfy legal obligations, and communicate with non-technical end-users (e.g., patients or bank customers). \textit{Red XAI} (model-validation oriented) serves as an active model interrogation framework for AI developers, researchers, and auditors to explore, validate, and debug black-box architectures.

\subsubsection{Project Alignment with Red XAI}\label{sec:red-xai-alignment}

This research is firmly rooted in the Red XAI paradigm by deploying evolutionary algorithms (EAs) as active model-interrogation tools rather than output-justification mechanisms. Traditional post-hoc attribution methods, such as SHAP and LIME, produce continuous, instance-level feature-importance scores that are mathematically convenient but difficult to audit for global biological coherence \citep{slack2020}. Furthermore, continuous attribution methods can produce visually convincing explanations even for overfitted models that have memorised noise.

In contrast, EAs extract discrete, symbolic \texttt{IF--THEN} rules that can be directly evaluated across competing criteria such as fidelity, accuracy, and comprehensibility \citep{yedjour2024} and, in this study, checked against biology. Rule extraction is fundamentally a multi-objective combinatorial optimisation problem that balances rule simplicity against faithfulness to the underlying model \citep{yedjour2024}, which makes evolutionary search a natural fit. Whether this approach scales to high-dimensional genomic feature spaces is part of what this study tests.